\documentclass[10pt,a4paper]{amsart}
\usepackage{amsmath,amssymb,amsthm,graphicx}
% Retained minimal preamble for the Stacks Project chapter varieties.tex.
% Verbatim excerpt of the chapter preamble: Xy-pic setup (native preamble lines 14-18), the native theorem
% declarations with the mathematical macros (native preamble lines 154-217), and the project own no-op
% definitions of the reference/slogan/history display environments (native wrapper.tex lines 1-3) that the
% counted statement uses for its slogan. Copied mechanically; nothing added, removed or redefined.
% For commutative diagrams we use Xy-pic
\usepackage[all]{xy}

% We use 2cell for 2-commutative diagrams.
\xyoption{2cell}
% Theorem environments.
%
\theoremstyle{plain}
\newtheorem{theorem}[subsection]{Theorem}
\newtheorem{proposition}[subsection]{Proposition}
\newtheorem{lemma}[subsection]{Lemma}

\theoremstyle{definition}
\newtheorem{definition}[subsection]{Definition}
\newtheorem{example}[subsection]{Example}
\newtheorem{exercise}[subsection]{Exercise}
\newtheorem{situation}[subsection]{Situation}

\theoremstyle{remark}
\newtheorem{remark}[subsection]{Remark}
\newtheorem{remarks}[subsection]{Remarks}

\numberwithin{equation}{subsection}

% Macros
%
\def\lim{\mathop{\mathrm{lim}}\nolimits}
\def\colim{\mathop{\mathrm{colim}}\nolimits}
\def\Spec{\mathop{\mathrm{Spec}}}
\def\Hom{\mathop{\mathrm{Hom}}\nolimits}
\def\Ext{\mathop{\mathrm{Ext}}\nolimits}
\def\SheafHom{\mathop{\mathcal{H}\!\mathit{om}}\nolimits}
\def\SheafExt{\mathop{\mathcal{E}\!\mathit{xt}}\nolimits}
\def\Sch{\mathit{Sch}}
\def\Mor{\mathop{\mathrm{Mor}}\nolimits}
\def\Ob{\mathop{\mathrm{Ob}}\nolimits}
\def\Sh{\mathop{\mathit{Sh}}\nolimits}
\def\NL{\mathop{N\!L}\nolimits}
\def\CH{\mathop{\mathrm{CH}}\nolimits}
\def\proetale{{pro\text{-}\acute{e}tale}}
\def\etale{{\acute{e}tale}}
\def\QCoh{\mathit{QCoh}}
\def\Ker{\mathop{\mathrm{Ker}}}
\def\Im{\mathop{\mathrm{Im}}}
\def\Coker{\mathop{\mathrm{Coker}}}
\def\Coim{\mathop{\mathrm{Coim}}}

% Boxtimes
%
\DeclareMathSymbol{\boxtimes}{\mathbin}{AMSa}{"02}

%
% Macros for moduli stacks/spaces
%
\def\QCohstack{\mathcal{QC}\!\mathit{oh}}
\def\Cohstack{\mathcal{C}\!\mathit{oh}}
\def\Spacesstack{\mathcal{S}\!\mathit{paces}}
\def\Quotfunctor{\mathrm{Quot}}
\def\Hilbfunctor{\mathrm{Hilb}}
\def\Curvesstack{\mathcal{C}\!\mathit{urves}}
\def\Polarizedstack{\mathcal{P}\!\mathit{olarized}}
\def\Complexesstack{\mathcal{C}\!\mathit{omplexes}}
% \Pic is the operator that assigns to X its picard group, usage \Pic(X)
% \Picardstack_{X/B} denotes the Picard stack of X over B
% \Picardfunctor_{X/B} denotes the Picard functor of X over B
\def\Pic{\mathop{\mathrm{Pic}}\nolimits}
\def\Picardstack{\mathcal{P}\!\mathit{ic}}
\def\Picardfunctor{\mathrm{Pic}}
\def\Deformationcategory{\mathcal{D}\!\mathit{ef}}
\newenvironment{reference}{}{}
\newenvironment{slogan}{}{}
\newenvironment{history}{}{}
\usepackage{cleveref}
\crefname{originallemma}{original lemma}{original lemmas}
\Crefname{originallemma}{Original lemma}{Original lemmas}
\makeatletter
\newlabel{lemma-open-in-normal-proper}{{04L4}{source}{}{sourcepointer-1}{}}
\newlabel{lemma-open-in-normal-proper@cref}{{[originallemma][0][]04L4}{[1][0][]source}{}{}{}}
\makeatother
\crefname{originaldefinition}{original definition}{original definitions}
\Crefname{originaldefinition}{Original definition}{Original definitions}
\makeatletter
\newlabel{morphisms-definition-normalization}{{035N}{source}{}{sourcepointer-2}{}}
\newlabel{morphisms-definition-normalization@cref}{{[originaldefinition][0][]035N}{[1][0][]source}{}{}{}}
\makeatother
\crefname{originallemma}{original lemma}{original lemmas}
\Crefname{originallemma}{Original lemma}{Original lemmas}
\makeatletter
\newlabel{morphisms-lemma-Japanese-normalization}{{035R}{source}{}{sourcepointer-3}{}}
\newlabel{morphisms-lemma-Japanese-normalization@cref}{{[originallemma][0][]035R}{[1][0][]source}{}{}{}}
\makeatother
\crefname{originallemma}{original lemma}{original lemmas}
\Crefname{originallemma}{Original lemma}{Original lemmas}
\makeatletter
\newlabel{morphisms-lemma-normalization-normal}{{035Q}{source}{}{sourcepointer-4}{}}
\newlabel{morphisms-lemma-normalization-normal@cref}{{[originallemma][0][]035Q}{[1][0][]source}{}{}{}}
\makeatother
\crefname{originallemma}{original lemma}{original lemmas}
\Crefname{originallemma}{Original lemma}{Original lemmas}
\makeatletter
\newlabel{morphisms-lemma-ubiquity-nagata}{{035B}{source}{}{sourcepointer-5}{}}
\newlabel{morphisms-lemma-ubiquity-nagata@cref}{{[originallemma][0][]035B}{[1][0][]source}{}{}{}}
\makeatother
\crefname{originalsection}{original section}{original sections}
\Crefname{originalsection}{Original section}{Original sections}
\makeatletter
\newlabel{morphisms-section-integral}{{01WG}{source}{}{sourcepointer-6}{}}
\newlabel{morphisms-section-integral@cref}{{[originalsection][0][]01WG}{[1][0][]source}{}{}{}}
\makeatother
\crefname{originalsection}{original section}{original sections}
\Crefname{originalsection}{Original section}{Original sections}
\makeatletter
\newlabel{morphisms-section-proper}{{01W0}{source}{}{sourcepointer-7}{}}
\newlabel{morphisms-section-proper@cref}{{[originalsection][0][]01W0}{[1][0][]source}{}{}{}}
\makeatother
\title{Finiteness of the group of units modulo constants for an integral locally finite type scheme over an algebraically closed field}
\author{The Stacks Project Authors}
\date{stacks-project:varieties.tex:local-snapshot-20261011; GNU FDL 1.2 or later}
\begin{document}
\maketitle
\noindent\textit{Editorial note.} All delivered mathematical text is the author's own, copied from the pinned Stacks Project chapter file varieties.tex; no statement or proof step was rewritten, reordered or removed.\par\medskip
\noindent\textit{Editorial note.} The proof is complete as written and ends by quoting the sibling normal proper lemma; the reduction that makes that lemma applicable is entirely inside the counted span.\par\medskip
\noindent\textit{Editorial note.} The counted spans carry 7 navigation stubs and no bibliography citation.\par\medskip
\noindent\textit{Editorial note.} The preamble is the chapter's retained minimal preamble (sources/stacks-varieties/retained-preamble.tex lines 1-77), a verbatim excerpt of the native preamble.\par\medskip
\noindent\textit{Editorial note.} The label lemma-units-integral-finite-type-algebraically-closed is defined only in varieties.tex, so no cross-chapter restatement had to be weighed for this card.\par\medskip
\noindent\textit{Allowed standard prerequisites (editorial).}\begin{itemize}
\item Normalization morphism: finite, dominant, with normal target. Morphisms, Definition morphisms-definition-normalization and Lemmas morphisms-lemma-normalization-normal, morphisms-lemma-Japanese-normalization, morphisms-lemma-ubiquity-nagata, invoked at proof lines 5264-5271. Supplies that the normalization used in the reduction is finite and that the normalized variety is normal.
\item Properness is stable under finite morphisms, base change and composition. Morphisms, Sections morphisms-section-proper and morphisms-section-integral, invoked at proof lines 5272-5276. Supplies that the normalization of a projective variety is again proper.
\item The normal proper case. Same chapter, Lemma lemma-open-in-normal-proper, invoked at proof lines 5280-5282. Supplies the finiteness statement once X is a nonempty open of a normal proper variety.
\end{itemize}\medskip
\ifdefined\hypertarget\hypertarget{sourcepointer-1}{}\fi
\noindent\textit{Source reference.} Original lemma 04L4: sources/stacks-varieties/pinned-original-source.tex, source lines 5157--5157. Pointer outside the counted span, invoked once for context; the load-bearing argument of the counted statement is the author's proof reproduced in full in the proof span.\par\medskip
\ifdefined\hypertarget\hypertarget{sourcepointer-2}{}\fi
\noindent\textit{Source reference.} Original definition 035N: sources/stacks-chow/tags.txt, source lines 4134--4134. Pointer outside the counted span, invoked once for context; the load-bearing argument of the counted statement is the author's proof reproduced in full in the proof span.\par\medskip
\ifdefined\hypertarget\hypertarget{sourcepointer-3}{}\fi
\noindent\textit{Source reference.} Original lemma 035R: sources/stacks-chow/tags.txt, source lines 4138--4138. Pointer outside the counted span, invoked once for context; the load-bearing argument of the counted statement is the author's proof reproduced in full in the proof span.\par\medskip
\ifdefined\hypertarget\hypertarget{sourcepointer-4}{}\fi
\noindent\textit{Source reference.} Original lemma 035Q: sources/stacks-chow/tags.txt, source lines 4137--4137. Pointer outside the counted span, invoked once for context; the load-bearing argument of the counted statement is the author's proof reproduced in full in the proof span.\par\medskip
\ifdefined\hypertarget\hypertarget{sourcepointer-5}{}\fi
\noindent\textit{Source reference.} Original lemma 035B: sources/stacks-chow/tags.txt, source lines 4122--4122. Pointer outside the counted span, invoked once for context; the load-bearing argument of the counted statement is the author's proof reproduced in full in the proof span.\par\medskip
\ifdefined\hypertarget\hypertarget{sourcepointer-6}{}\fi
\noindent\textit{Source reference.} Original section 01WG: sources/stacks-chow/tags.txt, source lines 2481--2481. Pointer outside the counted span, invoked once for context; the load-bearing argument of the counted statement is the author's proof reproduced in full in the proof span.\par\medskip
\ifdefined\hypertarget\hypertarget{sourcepointer-7}{}\fi
\noindent\textit{Source reference.} Original section 01W0: sources/stacks-chow/tags.txt, source lines 2465--2465. Pointer outside the counted span, invoked once for context; the load-bearing argument of the counted statement is the author's proof reproduced in full in the proof span.\par\medskip
\begin{lemma}
\label{lemma-units-integral-finite-type-algebraically-closed}
Let $k$ be an algebraically closed field.
Let $X$ be an integral scheme locally of finite type over $k$.
Then $\mathcal{O}^*(X)/k^*$ is a finitely generated abelian group.
\end{lemma}

\begin{proof}
As $X$ is integral the restriction mapping
$\mathcal{O}(X) \to \mathcal{O}(U)$ is injective for any
nonempty open subscheme $U \subset X$. Hence we may assume
that $X$ is affine. Choose a closed immersion
$X \to \mathbf{A}^n_k$
and denote $\overline{X}$ the closure of $X$ in $\mathbf{P}^n_k$
via the usual immersion $\mathbf{A}^n_k \to \mathbf{P}^n_k$.
Thus we may assume that $X$ is an affine open of a projective
variety $\overline{X}$.

\medskip\noindent
Let $\nu : \overline{X}^\nu \to \overline{X}$ be the normalization
morphism, see
Morphisms, Definition \ref{morphisms-definition-normalization}.
We know that $\nu$ is finite, dominant, and that $\overline{X}^\nu$
is a normal irreducible scheme, see
Morphisms, Lemmas \ref{morphisms-lemma-normalization-normal},
\ref{morphisms-lemma-Japanese-normalization}, and
\ref{morphisms-lemma-ubiquity-nagata}.
It follows that $\overline{X}^\nu$ is a proper variety,
because $\overline{X} \to \Spec(k)$ is proper as a composition
of a finite and a proper morphism (see results in
Morphisms, Sections \ref{morphisms-section-proper} and
\ref{morphisms-section-integral}).
It also follows that $\nu$ is a surjective morphism, because
the image of $\nu$ is closed and contains the generic point of $\overline{X}$.
Hence setting $X^\nu = \nu^{-1}(X)$ we see that it suffices to prove the
result for $X^\nu$. In other words, we may assume that $X$ is a nonempty
open of a normal proper variety $\overline{X}$. This case is handled by
Lemma \ref{lemma-open-in-normal-proper}.
\end{proof}

\end{document}
